\section{Problem Statement}
\label{sec:problem}

\subsection{Problem}

Network proximity can shortlist AD drug candidates, and complementary exposure tells us what a good pair looks like. Turning this into a recommended \emph{combination} of $k$ drugs is a search problem whose size grows as $\binom{n}{k}$. Three questions are open: how to turn closeness, overlap and combination size into one score without inventing biological numbers; how much one search algorithm gains from SIMD and from a GPU, and why; and how well a quantum algorithm (QAOA), simulated on a GPU, finds the best combination compared with the exact answer.

\subsection{Proposed solution}

We build the system as five modules. Each module has fixed inputs and outputs, so team members can build and test them separately (\Cref{tab:modules}). Modules 1--4 are the graded core; Module 5 is the research extension.

\begin{table}[H]
\centering
\small
\begin{tabularx}{\linewidth}{>{\raggedright\arraybackslash}p{3.0cm}>{\raggedright\arraybackslash}X>{\raggedright\arraybackslash}X>{\raggedright\arraybackslash}p{2.2cm}}
\toprule
Module & Input & Output & Runs on \\
\midrule
1. Candidate preparation & Interactome, drug--target list, AD genes, candidate list; \emph{or} the sister project's table & Candidate table (targets, $z_i$) and a shortlist of $n$ drugs & CPU (Python) \\
2. Problem construction & Shortlist, network distances & Rewards $w_i$, penalties $P_{ij}$, QUBO and Ising files & CPU (Python) \\
3. Exact search & $w$, $P$, $k$ & The best combination, how many combinations tie for best & CPU and GPU \\
4. Simulated annealing & $w$, $P$, $k$, settings & Best combination per run, run statistics, timings & CPU scalar, AVX2, GPU \\
5. QAOA (extension) & Ising file & Sampled combinations and their scores & GPU (CUDA-Q + cuStateVec) \\
\bottomrule
\end{tabularx}
\caption{Pipeline modules. Every solver reads the same problem file and writes the same result format (\Cref{app:data}).}
\label{tab:modules}
\end{table}

\subsection{Research questions and hypotheses}

\begin{description}[leftmargin=0.8cm,style=nextline]
\item[RQ1 (main): How much do SIMD and the GPU help?] How do sequential C++, AVX2 and CUDA versions of the \emph{same} simulated-annealing algorithm compare in speed, as the number of candidates $n$ grows? We measure the time for one search chain, the number of steps per second when many chains run at once, and the time until the exact best answer is found.\\
\emph{Hypothesis:} The GPU wins once enough independent chains run to keep it busy (hundreds or more). AVX2 wins for a single chain. Inside a chain, the speedup is limited because only \emph{accepted} steps do the vectorized array update, and late in the search most steps are rejected.
\item[RQ2 (main): Is the scoring function right?] Does our QUBO always pick a valid combination of exactly $k$ drugs, equal to the best combination found by exact search, and how does the overlap dial $\lambda$ change which drugs are chosen?\\
\emph{Hypothesis:} With the penalty weight set above the bound we prove in \Cref{app:formulation}, the QUBO and exact answers always agree. As $\lambda$ grows, the chosen combinations move from ``the $k$ drugs closest to AD'' toward complementary combinations.
\item[RQ3 (extension): How good is QAOA?] On instances small enough to solve exactly, how often does simulated QAOA return the best combination, how many of its answers are valid $k$-drug combinations, and how does this change with circuit depth $p$ and problem size $n$?\\
\emph{Hypothesis:} Most of QAOA's answers will be invalid, because only a tiny share of all yes/no patterns have exactly $k$ ones ($\binom{n}{k}/2^n$; 0.1\% at $n=20$, $k=3$). Quality will improve with $p$ but remain well below exact search.
\item[RQ4 (supporting): Where does the time go?] How is total run time split between data preparation, problem construction, solving and checking? Answered from timing logs at no extra cost.
\end{description}

RQ1 and RQ2 are the course core and do not depend on the quantum part. RQ3 is the main research question for a possible conference paper.

\subsection{Objectives}

\textbf{General objective.} To build a checked scoring function for AD drug combinations and compare exact, SIMD, GPU and simulated-quantum ways of finding the best combination.

\textbf{Specific objectives.}
\begin{enumerate}
\item Build a reproducible candidate table and shortlist from public data, recording where every number comes from.
\item Build the QUBO and Ising files and prove, then test, that they give the right answer.
\item Implement exact search on CPU and GPU and confirm they agree.
\item Implement one simulated-annealing algorithm in sequential C++, AVX2 assembly and CUDA; show that all three make identical decisions; compare their execution times.
\item Run QAOA in CUDA-Q on the GPU simulator and measure its answers against the exact search.
\end{enumerate}

\subsection{Scope and limitations}

\begin{itemize}
\item \textbf{Public data only.} No patient data.
\item \textbf{Hypotheses, not recommendations.} A high-scoring combination is a computer-generated idea. It is not evidence that the combination is effective, synergistic or safe.
\item \textbf{One mid-range GPU.} A GeForce GTX 1660 with 6 GB. This limits the quantum simulation to about 28 qubits (one qubit per drug). It does not limit the classical solvers at our sizes.
\item \textbf{No quantum speedup is claimed.} The quantum computer is simulated, and at these sizes exact search is fast. The quantum part measures answer quality only.
\item \textbf{Weak biological validation.} Very few AD drug combinations have trial results, so we check whether chosen combinations follow the complementary-exposure pattern, not whether they work.
\item \textbf{Cut for time.} Safety and side-effect data, splitting large problems into QAOA-sized pieces, a constraint-preserving QAOA variant, and several extra performance experiments were left out. \Cref{app:limitations} lists each one, why it was cut, and how it would be added later.
\end{itemize}

\subsection{Significance}

For the course, the project shows one algorithm on three kinds of hardware and explains where data-level parallelism helps and where a step-by-step dependency limits it. For research, it gives a realistic, biology-based test problem for quantum optimization, with exact answers to compare against, and the first QAOA study of AD drug combinations that we are aware of.
